\hmsubsection{Similar Systems}{FS}{}\label{sec:related-systems}
One earlier bachelor project at USN addressed sorting tasks that are
closely related to this project. Industrial pick-and-place cells provide
a third reference point. Considering these systems together clarifies
which design elements were reused, adapted, or replaced, and places this
project on the scale between a small academic demonstrator and an
industrial system.

\subsubsection*{Sortify (USN, 2025)}
Sortify is the closest comparator among the reviewed student projects
\cite{sortify}. That system sorted red and blue balls using a Raspberry
Pi 5 host, an OAK-D camera, and a Dynamixel-based arm. Its vision
pipeline used CLAHE, HSV segmentation, the Hough Circle Transform, and a
YOLOv4-tiny detector as an additional verifier. The project also
implemented a custom UART protocol with CRC-16 framing.
The same CLAHE--HSV--Hough core was retained in this project
(Section \ref{sec:image-processing}) because it is well suited to the
sorting task and already satisfies the relevant detection requirements.
However, three changes were introduced.
First, no additional YOLO model was added. Instead, HSV and Hough
detections are merged using a Union-Find pass and a Non-Maximum
Suppression step (Section \ref{sec:ensemble}). This ensemble serves a
similar verification role while avoiding the cost of training and
running a CNN.
Second, the wire protocol uses plain text. Sortify used CRC-16 because
the arm was intended for installation on a moving vehicle, where
electrical noise is a more significant concern. In this project, the
system operates on a bench with cable lengths below one metre. Text
framing simplifies terminal-based debugging, and the available
development observations did not indicate transmission issues. CRC-16 can
be added if the arm is later mounted on the Fenris vehicle.
Third, Sortify reported in its conclusion that its inverse kinematics
was not accurate enough to complete fully autonomous pick-and-place. This
limitation influenced the present design. The solver used here is a
closed-form geometric solver (Section \ref{sec:ik-solver}) rather than
the iterative method used by Sortify. The solver is also paired with a
sag compensation step (Section \ref{sec:sag}) and per-point surface
heights from touch calibration (Section \ref{sec:touch-cal}).
\subsubsection*{Industrial pick-and-place cells}
At the other end of the scale, industrial systems include cells such as
the Universal Robots UR3e with a wrist-mounted camera, as well as
dedicated sorters from Tomra. These systems typically use six- or
seven-axis arms, structured lighting, and detection models trained on
application-specific datasets. They are also substantially more
expensive than a bachelor project and require formal installation work.
This project is not intended to compete directly with industrial
systems. Instead, it occupies a position between the USN bachelor
projects discussed above and commercial pick-and-place cells. The goal
is to produce a functional demonstrator that can be explained end to
end, rather than the fastest or most flexible sorter available.
